\documentclass[10pt,a4paper]{amsart}
\usepackage[margin=19mm]{geometry}
\usepackage{amsmath,amssymb,mathtools}
\usepackage[scr=rsfs,cal=euler]{mathalfa}
\usepackage{xfrac}
\usepackage[unicode,hidelinks,bookmarks=false]{hyperref}
\newcommand{\ie}{i.e.\ }
\newcommand{\cf}{cf.\ }
\newcommand{\resp}{respectively\ }
\newcommand{\Subsection}{\subsection}
\providecommand{\nobreakdash}{\nobreak\hbox{-}}
\setlength{\emergencystretch}{2em}
\multlinegap0pt
\usepackage{bm}
\usepackage{bbm}
\usepackage{dsfont}
\usepackage{xspace}
\usepackage{cancel}
\usepackage{mathtools}
\usepackage{amsthm}
\makeatletter
\@ifundefined{thm}{\newtheorem{thm}{Thm}}{}
\@ifundefined{lem}{\newtheorem{lem}[thm]{Lemma}}{}
\makeatother
\makeatletter
\expandafter\ifx\csname customlem\endcsname\relax
\newenvironment{customlem}[1]
{\renewcommand\theinnercustomlem{#1}\innercustomlem}
{\endinnercustomlem}
\fi
\expandafter\ifx\csname customprop\endcsname\relax
\newenvironment{customprop}[1]
{\renewcommand\theinnercustomprop{#1}\innercustomprop}
{\endinnercustomprop}
\fi
\expandafter\ifx\csname customthm\endcsname\relax
\newenvironment{customthm}[1]
{\renewcommand\theinnercustomthm{#1}\innercustomthm}
{\endinnercustomthm}
\fi
\newcommand*{\defeq}{\mathrel{\rlap{%
			\raisebox{0.3ex}{$\m@th\cdot$}}%
		\raisebox{-0.3ex}{$\m@th\cdot$}}%
	=}
\renewcommand{\phi}{\varphi}
\makeatother
\title[The fully nonlinear Loewner-Nirenberg problem: Liouville the]{The fully nonlinear Loewner-Nirenberg problem: Liouville theorems and counterexamples to local boundary estimates}
\author{Jonah A. J. Duncan, Luc Nguyen}
\date{Source: arXiv:2507.16383v1 (CC BY 4.0)}
\begin{document}
\maketitle

\noindent\emph{Source and licence.} The fully nonlinear Loewner-Nirenberg problem: Liouville theorems and counterexamples to local boundary estimates. Version arXiv:2507.16383v1 (CC BY 4.0), distributed under the Creative Commons Attribution 4.0 International licence, as recorded in the arXiv licence metadata for this exact version and shown on the abstract page https://arxiv.org/abs/2507.16383v1. Full rights notice: licenses/arxiv-2507.16383-CC-BY-4.0.txt.\par\smallskip

\noindent\textit{Editorial scope.} Counted unit: the Lem whose source label is \texttt{l3} in the article (source0.tex lines 472--489, source label \texttt{l3}), delivered together with the article's own complete original proof of it (source0.tex lines 491--528, one \texttt{proof} environment of 38 lines). No mathematics is written, repaired or re-derived: statement and proof are the article's own text line for line. Required context retained uncounted: none. Every definition, display and label of those lines is retained unchanged. Omitted from this package and not redistributed: 1--471, 490--490, 529--1506. Nothing inside a retained excerpt is omitted or altered: every retained body is delivered exactly as the article's lines are written, so no omission receipt of any kind accompanies this package. Citations: the retained excerpts carry no citation command at all, so no bibliography entry of the article is transcribed and no reference is redistributed. Reference adaptation: none was needed and none was made; every reference inside the delivered unit resolves inside it, so no unresolved reference or citation is produced. Printed slips kept literal: no printed word, symbol, hypothesis or number of the counted statement or of its proof was altered. Layout and class: the article's own class is \texttt{article} with packages including amsmath, enumerate, amsfonts, amssymb, color, graphicx, amsthm, hyperref, todonotes, ulem, cite, mathtools, tocbibind, appendix; the wrapper is the lane's pinned \texttt{amsart} editorial wrapper at 10pt on A4, which restates the theorem-like environments the retained text needs and adds the article's own macro definitions that the retained text needs (\texttt{\textbackslash defeq}, \texttt{\textbackslash phi}), with extra wrapper packages (bm, bbm, dsfont, xspace, cancel, mathtools). The delivered numbering is the unit's own. Assets: no retained span contains a figure environment, a graphics-inclusion command, a PDF-page inclusion, a table or a TikZ picture; no image file is copied into this package, the package declares no asset dependency and no image file is redistributed with it. Licence: version arXiv:2507.16383v1 of this article is distributed under the Creative Commons Attribution 4.0 International licence, as recorded in the arXiv licence metadata for this exact version and shown on the abstract page https://arxiv.org/abs/2507.16383v1. A full rights notice is delivered at licenses/arxiv-2507.16383-CC-BY-4.0.txt. Characters: every retained excerpt is pure ASCII, so the delivered engine, XeLaTeX with its default Latin Modern text font, reports no missing character.
\begin{lem}\label{l3}
	Fix $x_0'\in\{x_n=0\}$ and suppose $R>200$, $\frac{r_1^2}{R^2} < \frac{1}{200}$ and $ \sqrt{R^2 - r_1^2} \leq r < R$. Let $A_{r, \sqrt{R^2 + r_1^2}}$ denote the open annulus centred at $(x_0', -R)$ with inner radius $r$ and outer radius $\sqrt{R^2 + r_1^2}$, and for $x = (x',x_n)$ let $d(x) = \sqrt{|x' - x_0'|^2 + (x_n + R)^2}$ denote the distance to its centre. Then for any constant $C$ satisfying
	\begin{align}\label{11}
	\frac{9R^2}{r_1^4} < C < \frac{R^3}{21r_1^4}, 
	\end{align}
	the conformal metric $g_{h_{r}}$ defined on $A_{r, \sqrt{R^2 + r_1^2}}$ by
	\begin{align*}
	h_{r}(x) = h_{r}(d) = (d-r) + C(d-r)^2 > 0 
	\end{align*}
	satisfies
	\begin{align*}
	\begin{cases}
	f(\lambda(-A_{h_{r}})) \geq \frac{1}{2}, \quad \lambda(-A_{h_r})\in\Gamma & \text{in }A_{r, \sqrt{R^2 + r_1^2}} \\
	h_{r} = 0 & \text{on }\{d=r\} \\
	h_{r} \geq 1 & \text{on }\{d=\sqrt{R^2 + r_1^2}\}. 
	\end{cases}
	\end{align*}
\end{lem}

\begin{proof}
	The fact that $h_r = 0$ on $\{d=r\}$ is clear from the definition of $h_r$. To see that $h_r \geq 1$ on $\{d=\sqrt{R^2+r_1^2}\}$, it suffices to show that $C(\sqrt{R^2 + r_1^2} - R)^2 \geq 1$ under our hypotheses on $R, r_1, r$ and $C$. To this end, note that since $\sqrt{1+x} \geq 1 + \frac{x}{3}$ for $x\in[0,3]$ and since we assume $\frac{r_1^2}{R^2}< \frac{1}{200}$, we have the inequality
	\begin{align}\label{12}
	\sqrt{R^2 + r_1^2} = R\bigg(1+ \frac{r_1^2}{R^2}\bigg)^{1/2} \geq R\bigg(1+\frac{r_1^2}{3R^2}\bigg).
	\end{align}
	The inequality in \eqref{12} combined with the lower bound on $C$ in \eqref{11} then clearly implies $C(\sqrt{R^2 + r_1^2} - R)^2 \geq 1$, as required.\medskip 
	
	
	
	To complete the proof of the claim, it remains to show that our assumptions on $R, r_1, r$ and $C$ imply $f(\lambda(-A_{h_{r}})) \geq \frac{1}{2}$ and $\lambda(-A_{h_{r}})\in\Gamma$ in $A_{r, \sqrt{R^2 + r_1^2}}$. We start by computing the eigenvalues of $-A_{h_{r}}$, which have multiplicity 1 and $(n-1)$, respectively: 
	\begin{align*}
	\lambda_1 & = -h_{r}h_{r}'' + \frac{1}{2}(h_{r}')^2 = -2C\big((d-r) + C(d-r)^2\big) + \frac{1}{2}\big(1+2C(d-r)\big)^2 = \frac{1}{2}
	\end{align*}
	and
	\begin{align*}
	\lambda_2 & = -\frac{h_{r}h_{r}'}{d} + \frac{1}{2}(h_{r}')^2   = -\big((d-r) + C(d-r)^2\big)\frac{1+2C(d-r)}{r}+ \frac{1}{2}\big(1+2C(d-r)\big)^2 \nonumber \\
	& = \frac{1}{2} + (d-r)\bigg[2C -\frac{1}{d} + \Big(\frac{-3C}{d} + 2C^2\Big)(d-r) - \frac{2C^2}{d}(d-r)^2\bigg] \nonumber \\
	& \geq \frac{1}{2} + (d-r)\bigg[2C -\frac{1}{d} -\frac{3C}{d}(d-r) - \frac{2C^2}{d}(d-r)^2\bigg] \nonumber \\
	& \geq \frac{1}{2} + (d-r)\bigg[2C -\frac{7}{4d} - \frac{5C^2}{d}(d-r)^2\bigg],
	\end{align*}
	where we have applied the Cauchy-Schwarz inequality to reach the last line. It therefore suffices to show that
	\begin{align}\label{500}
	\phi \defeq 2C -\frac{7}{4d} - \frac{5C^2}{d}(d-r)^2  \geq 0. 
	\end{align} 
	To prove \eqref{500} we first note
	\begin{align*}
	R\bigg(1-\frac{r_1^2}{R^2}\bigg) \leq \sqrt{R^2 - r_1^2}\, \leq r < d < \sqrt{R^2 + r_1^2} \, \leq R\bigg(1+\frac{r_1^2}{R^2}\bigg),
	\end{align*} 
	which implies $d-r \leq \frac{2r_1^2}{R}$. These, together with $\frac{r_1^2}{R^2} < \frac{1}{200}$, imply
	\begin{align*}
	\frac{7}{4d} < \frac{7}{4R(1-\frac{r_1^2}{R^2})} < \frac{2}{R} \quad \text{and}\quad 	\frac{5C^2}{d}(d-r)^2 < \frac{5C^2}{R(1-\frac{r_1^2}{R^2})}\frac{4r_1^4}{R^2} < \frac{21C^2 r_1^4}{R^3}. 
	\end{align*}
	Therefore 
	\begin{align*}
	\phi \geq 2C - \frac{2}{R} - \frac{21C^2 r_1^4}{R^3},
	\end{align*}
	and thus $\phi$ is clearly seen to be nonnegative by \eqref{11}. 
\end{proof}

\end{document}
